\renewcommand{\bibname}{参考资料与实验出处}
\begin{thebibliography}{99}
    \addcontentsline{toc}{chapter}{参考资料与实验出处}
    \item[] 正文中的参考条目按教学主题排列。教材与讲义用于进一步学习推导，实验论文及作者回顾提供装置、观测和历史背景。
    一维波动力学、谐振子与氢原子的部分推导整理自作者2021年的《普通物理 II》课堂笔记，授课教师为 Shuo Jiang。

    \bibitem{mitwave}
    B. Zwiebach，MIT 8.04 Quantum Physics I，Spring 2016。讲义5--7讨论波函数、概率流与自由波包，讲义11讨论方势阱，讲义16、17讨论势阶散射与透射共振。
    \href{https://ocw.mit.edu/courses/8-04-quantum-physics-i-spring-2016/pages/lecture-notes/}{课程讲义索引}。
    \bibitem{mitspin}
    B. Zwiebach，MIT 8.05 Quantum Physics II，Fall 2013，讲义2：
    \textit{Spin One-half, Bras, Kets, and Operators}。
    \href{https://ocw.mit.edu/courses/8-05-quantum-physics-ii-fall-2013/resources/lecture-notes-2/}{公开讲义}。
    \bibitem{mituncertainty}
    B. Zwiebach，同课程，讲义5：\textit{Uncertainty Principle and Compatible Observables}。
    \href{https://ocw.mit.edu/courses/8-05-quantum-physics-ii-fall-2013/resources/lecture-notes-5/}{公开讲义}。
    \bibitem{mittwostate}
    B. Zwiebach，同课程，讲义7：\textit{Two State Systems}。
    \href{https://ocw.mit.edu/courses/8-05-quantum-physics-ii-fall-2013/resources/lecture-notes-7/}{公开讲义}。
    \bibitem{mitangular}
    B. Zwiebach，同课程，讲义9、10：\textit{Angular Momentum; Addition of Angular Momentum}。
    \href{https://ocw.mit.edu/courses/8-05-quantum-physics-ii-fall-2013/pages/lecture-notes/}{课程讲义索引}。
    \bibitem{mitpert}
    B. Zwiebach，MIT 8.06 Quantum Physics III，Spring 2018，讲义1：
    \textit{Non-degenerate and Degenerate Perturbation Theory}，尤其1.1与1.2节。
    \href{https://ocw.mit.edu/courses/8-06-quantum-physics-iii-spring-2018/resources/chapter-1-non-degenerate-and-degenerate-perturbation-theory/}{公开讲义}。
    \bibitem{mitfine}
    B. Zwiebach，同课程，讲义2：\textit{Hydrogen Fine Structure}，尤其2.4节的相对论、旋轨与Darwin修正。
    \href{https://ocw.mit.edu/courses/8-06-quantum-physics-iii-spring-2018/resources/chapter-2-hydrogen-fine-structure/}{公开讲义}。
    \bibitem{mittime}
    B. Zwiebach，同课程，讲义4：\textit{Time Dependent Perturbation Theory}，尤其4.3节的黄金规则。
    \href{https://ocw.mit.edu/courses/8-06-quantum-physics-iii-spring-2018/resources/chapter-4-time-dependent-perturbation-theory/}{公开讲义}。
    \bibitem{mitadiabatic}
    B. Zwiebach，同课程，讲义6：\textit{Adiabatic Approximation}，绝热展开与Berry相位。
    \href{https://ocw.mit.edu/courses/8-06-quantum-physics-iii-spring-2018/resources/chapter-6-adiabatic-approximation/}{公开讲义}。
    \bibitem{mitscattering}
    B. Zwiebach，同课程，讲义7：\textit{Scattering}，尤其7.3节的积分方程和Born近似。
    \href{https://ocw.mit.edu/courses/8-06-quantum-physics-iii-spring-2018/resources/chapter-7-scattering/}{公开讲义}。
    \bibitem{lz}
    B. Zwiebach，同课程，L16.4--L16.5：\textit{Landau--Zener Transitions}。
    \href{https://ocw.mit.edu/courses/8-06-quantum-physics-iii-spring-2018/resources/l16-4/}{课程视频与转录入口}。
    \bibitem{feynman4}
    R. P. Feynman, R. B. Leighton, M. Sands，
    \textit{The Feynman Lectures on Physics}，Vol. III，第4章，Caltech在线版。
    \href{https://www.feynmanlectures.caltech.edu/III_04.html}{Identical Particles}。
    \bibitem{feynman9}
    同书，第9章。
    \href{https://www.feynmanlectures.caltech.edu/III_09.html}{The Ammonia Maser}。
    \bibitem{feynman18}
    同书，第18章。
    \href{https://www.feynmanlectures.caltech.edu/III_18.html}{Angular Momentum}。
    \bibitem{domains}
    G. Bonneau, J. Faraut, G. Valent，
    \textit{Self-adjoint extensions of operators and the teaching of quantum mechanics}，
    Am. J. Phys. \textbf{69}, 322--331 (2001)，尤其第5节及附录A。
    \href{https://arxiv.org/abs/quant-ph/0103153}{作者预印本}。
    \bibitem{stone}
    M. H. Stone，\textit{On One-Parameter Unitary Groups in Hilbert Space}，
    Annals of Mathematics \textbf{33} (1932), 643--648。
    \href{https://doi.org/10.2307/1968538}{原论文}。
    \bibitem{runge}
    B. Zwiebach，MIT 8.05，2013，作业11第6页（Runge--Lenz守恒量的逐项检验）；
    P. Etingof，\textit{Lie Groups and Lie Algebras}，氢原子隐藏对称性习题。
    \href{https://ocw.mit.edu/courses/8-05-quantum-physics-ii-fall-2013/33244de5a887410f525661ddf08e7103_MIT8_05F13_ps11.pdf}{守恒量计算}；
    \href{https://math.mit.edu/~etingof/lnlg.pdf}{Lie代数讲义}。
    \bibitem{manzano}
    D. Manzano，\textit{A short introduction to the Lindblad master equation}，
    AIP Advances \textbf{10}, 025106 (2020)。
    \href{https://arxiv.org/abs/1906.04478}{作者预印本}。

    \bibitem{littlejohnwe}
    R. G. Littlejohn，Physics 221A，Notes 20，
    \textit{The Wigner--Eckart Theorem} (2021)。
    \href{https://bohr.physics.berkeley.edu/classes/221/notes/wigeck.pdf}{Berkeley公开讲义}。

    \bibitem{littlejohnfw}
    R. G. Littlejohn，Physics 221B，Notes 49，
    \textit{The Foldy--Wouthuysen Transformation} (2021)，第4--8节。
    \href{https://bohr.physics.berkeley.edu/classes/221/notes/foldyw.pdf}{Berkeley公开讲义}。

    \bibitem{history}
    M. Planck，\textit{The Genesis and Present State of Development of the Quantum Theory}，
    Nobel Lecture (1920)，热辐射实验与量子假说的历史回顾；
    P. Lenard，Nobel Lecture (1906)，关于1887年Hertz光电现象的记录。
    \href{https://www.nobelprize.org/prizes/physics/1918/planck/lecture/}{Planck演讲}；
    \href{https://www.nobelprize.org/uploads/2017/07/lenard-lecture2.pdf}{Lenard演讲}。
    Balmer年代参见\href{https://www.nobelprize.org/uploads/2018/06/lamb-lecture.pdf}{Lamb的1955年Nobel演讲}；
    Zeeman年代参见\href{https://www.nobelprize.org/prizes/themes/the-nobel-prize-in-physics-1901-2000/}{Nobel物理学史资料}。
    光电定量检验见R. A. Millikan，Phys. Rev. \textbf{7}, 355 (1916)，
    \href{https://doi.org/10.1103/PhysRev.7.355}{A Direct Photoelectric Determination of Planck's $h$}。
    \bibitem{zeemanlecture}
    P. Zeeman，\textit{Light Radiation in a Magnetic Field}，Nobel Lecture (1902)；
    H. A. Lorentz，\textit{The Theory of Electrons and the Propagation of Light}，同年演讲。
    \href{https://www.nobelprize.org/prizes/physics/1902/zeeman/lecture/}{Zeeman演讲}；
    \href{https://www.nobelprize.org/prizes/physics/1902/lorentz/lecture/}{Lorentz演讲}。

    \bibitem{berryneutron}
    T. Bitter, D. Dubbers，
    \textit{Manifestation of Berry's Topological Phase in Neutron Spin Rotation}，
    Phys. Rev. Lett. \textbf{59}, 251--254 (1987)。
    \href{https://doi.org/10.1103/PhysRevLett.59.251}{原论文}。

    \bibitem{sg}
    W. Gerlach, O. Stern，\textit{Der experimentelle Nachweis der Richtungsquantelung im Magnetfeld}，
    Z. Phys. \textbf{9}, 349--352 (1922)。
    \href{https://doi.org/10.1007/BF01326983}{原论文}。
    历史解释的辨析另见B. Friedrich, D. Herschbach，
    \href{https://arxiv.org/abs/1805.09412}{The Stern--Gerlach Experiment and the Electron Spin}。
    \bibitem{compton}
    A. H. Compton，\textit{A Quantum Theory of the Scattering of X-rays by Light Elements}，
    Phys. Rev. \textbf{21}, 483 (1923)。
    \href{https://doi.org/10.1103/PhysRev.21.483}{原论文}。
    \bibitem{davisson}
    C. Davisson, L. H. Germer，\textit{Diffraction of Electrons by a Crystal of Nickel}，
    Phys. Rev. \textbf{30}, 705 (1927)。
    \href{https://doi.org/10.1103/PhysRev.30.705}{原论文}。
    \bibitem{franck}
    J. Franck，Nobel Lecture (1926)，\textit{Transformations of Kinetic Energy of Free Electrons into Excitation Energy of Atoms by Impacts}，
    回顾1914年Franck--Hertz实验。
    \href{https://www.nobelprize.org/uploads/2018/06/franck-lecture.pdf}{演讲入口}。
    \bibitem{rabi}
    I. I. Rabi, J. R. Zacharias, S. Millman, P. Kusch，
    \textit{A New Method of Measuring Nuclear Magnetic Moment}，
    Phys. Rev. \textbf{53}, 318 (1938)。
    \href{https://doi.org/10.1103/PhysRev.53.318}{原论文}。
    \bibitem{ramsey}
    N. F. Ramsey，\textit{A Molecular Beam Resonance Method with Separated Oscillating Fields}，
    Phys. Rev. \textbf{78}, 695 (1950)。
    \href{https://doi.org/10.1103/PhysRev.78.695}{方法论文}；1949年提出该方法的回顾见
    \href{https://www.nobelprize.org/uploads/2018/06/ramsey-lecture.pdf}{Ramsey的1989年Nobel演讲}。
    \bibitem{lamb}
    W. E. Lamb, Jr., R. C. Retherford，
    \textit{Fine Structure of the Hydrogen Atom by a Microwave Method}，
    Phys. Rev. \textbf{72}, 241 (1947)。
    \href{https://doi.org/10.1103/PhysRev.72.241}{原论文}。
    \bibitem{abexp}
    R. G. Chambers，\textit{Shift of an Electron Interference Pattern by Enclosed Magnetic Flux}，
    Phys. Rev. Lett. \textbf{5}, 3 (1960)。
    \href{https://doi.org/10.1103/PhysRevLett.5.3}{原论文}。
    \bibitem{qhe}
    K. v. Klitzing, G. Dorda, M. Pepper，
    \textit{New Method for High-Accuracy Determination of the Fine-Structure Constant Based on Quantized Hall Resistance}，
    Phys. Rev. Lett. \textbf{45}, 494 (1980)。
    \href{https://doi.org/10.1103/PhysRevLett.45.494}{原论文}。
    \bibitem{grangier}
    P. Grangier, G. Roger, A. Aspect，
    \textit{Experimental Evidence for a Photon Anticorrelation Effect on a Beam Splitter: A New Light on Single-Photon Interferences}，
    Europhys. Lett. \textbf{1}, 173--179 (1986)。
    \href{https://doi.org/10.1209/0295-5075/1/4/004}{原论文}。
    \bibitem{hom}
    C. K. Hong, Z. Y. Ou, L. Mandel，
    \textit{Measurement of Subpicosecond Time Intervals between Two Photons by Interference}，
    Phys. Rev. Lett. \textbf{59}, 2044 (1987)。
    \href{https://doi.org/10.1103/PhysRevLett.59.2044}{原论文}。
    \bibitem{tonomura}
    A. Tonomura等，\textit{Demonstration of Single-Electron Buildup of an Interference Pattern}，
    Am. J. Phys. \textbf{57}, 117--120 (1989)。
    \href{https://doi.org/10.1119/1.16104}{原论文}。
    \bibitem{belltests}
    S. J. Freedman, J. F. Clauser，Phys. Rev. Lett. \textbf{28}, 938 (1972)，
    \href{https://doi.org/10.1103/PhysRevLett.28.938}{早期光子检验}；
    A. Aspect, J. Dalibard, G. Roger，Phys. Rev. Lett. \textbf{49}, 1804 (1982)，
    \href{https://doi.org/10.1103/PhysRevLett.49.1804}{随时间切换分析器}；
    B. Hensen等，Nature \textbf{526}, 682--686 (2015)，
    \href{https://doi.org/10.1038/nature15759}{同时关闭主要漏洞的电子自旋检验}。
\end{thebibliography}
